% SPDX-License-Identifier: MIT
%
% structure.tex - Structural settings for the APA7 paper template.
%
% Controls how the PDF is built and laid out rather than what it says: language
% and academic period defaults, academic period output, and localized labels and
% section titles. Document values such as title, authors, and advisor live in
% data/metadata.tex.
%
% Language and academic period are resolved in this order, highest priority
% first:
%
%   1. compile.sh flags        ./compile.sh --language spanish --period semester
%   2. Defaults in this file   \documentLanguageDefault, \academicPeriodDefault
%   3. Built-in defaults       english, year
%
% compile.sh passes the final values to TeX, so compile through it rather
% than calling pdflatex directly.

\def\documentLanguageDefault{english}
\def\academicPeriodDefault{year}

\edef\documentLanguage{\documentLanguageDefault}
\edef\academicPeriod{\academicPeriodDefault}
\edef\metadataLanguageOverride{\documentLanguageOverride}
\edef\metadataPeriodOverride{\academicPeriodOverride}

\ifx\metadataLanguageOverride\empty\else
  \edef\documentLanguage{\metadataLanguageOverride}
\fi

\ifx\metadataPeriodOverride\empty\else
  \edef\academicPeriod{\metadataPeriodOverride}
\fi

%
% Academic period
%

% academic_period - Select year or semester output.
%
% The boolean is derived from academicPeriod so manual configuration and
% compile-time flags use the same code path.
\newif\ifusesemester
\edef\metadataSelectedPeriod{\academicPeriod}%
\def\metadataSemesterPeriod{semester}%
\ifx\metadataSelectedPeriod\metadataSemesterPeriod
  \usesemestertrue
\else
  \usesemesterfalse
\fi

\newcommand{\academicYear}{\the\year-\the\numexpr\year+1\relax}
\newcommand{\academicSemester}{\the\year-\ifnum\month<7 S01\else S02\fi}

%
% Localization
%

% localized_labels - Select labels from the resolved language.
%
% Unknown language values fall back to English. compile.sh validates known values
% before TeX starts, but this fallback keeps direct TeX runs predictable.
%
% A new section added to main.tex needs its own \sectionTitle... entry in both
% branches below.
\edef\metadataSelectedLanguage{\documentLanguage}%
\def\metadataSpanishLanguage{spanish}%
\ifx\metadataSelectedLanguage\metadataSpanishLanguage
  \def\labelAuthors{Autores}
  \def\labelAdvisor{Orientador}
  \def\labelAcademicYear{Año Académico}
  \def\labelAcademicSemester{Semestre Académico}
  \def\labelDate{Fecha}

  \def\sectionTitleAbstract{Resumen}
  \def\sectionTitleKeywords{Palabras clave}
  \def\sectionTitleAcknowledgments{Agradecimientos}
  \def\sectionTitleIntroduction{Introducción}
  \def\sectionTitleTheoreticalFramework{Marco teórico}
  \def\sectionTitleMethodology{Metodología}
  \def\sectionTitleResults{Resultados}
  \def\sectionTitleAnalysis{Análisis}
  \def\sectionTitleConclusion{Conclusión}
  \def\sectionTitleGlossary{Glosario}
  \def\sectionTitleReferences{Referencias}
  \def\sectionTitleAppendices{Apéndices}
\else
  \def\labelAuthors{Authors}
  \def\labelAdvisor{Advisor}
  \def\labelAcademicYear{Academic Year}
  \def\labelAcademicSemester{Academic Semester}
  \def\labelDate{Date}

  \def\sectionTitleAbstract{Abstract}
  \def\sectionTitleKeywords{Keywords}
  \def\sectionTitleAcknowledgments{Acknowledgments}
  \def\sectionTitleIntroduction{Introduction}
  \def\sectionTitleTheoreticalFramework{Theoretical Framework}
  \def\sectionTitleMethodology{Methodology}
  \def\sectionTitleResults{Results}
  \def\sectionTitleAnalysis{Analysis}
  \def\sectionTitleConclusion{Conclusion}
  \def\sectionTitleGlossary{Glossary}
  \def\sectionTitleReferences{References}
  \def\sectionTitleAppendices{Appendices}
\fi
